\section*{Impact Statement}

This work exposes a critical vulnerability in the measurement of AI progress: the ``illusion of competence'' arising from test set contamination. We demonstrate that even minimal data leakage allows models to bypass fundamental scaling laws through fragile, verbatim memorization rather than genuine reasoning. This creates a significant societal risk where models may be deployed based on inflated metrics, only to fail brittlely in novel, real-world scenarios where memorization offers no advantage.

Ethically, our findings mandate a shift toward more rigorous data curation and forensic evaluation practices. By identifying specific ``levers'' that expose memorization, such as sampling temperature and solution length, we provide actionable tools to restore trust in generative benchmarks. As AI systems become increasingly integrated into society, ensuring that evaluations reflect robust generalization rather than rote recall is paramount for safety and reliability.